\section{The Slack Principle}\label{sec:slack}

\Cref{lem:bidding} and \cref{prop:master} price the chokepoint. On our reading of the record since 2019 (\cref{sec:evidence}), Europe's premium has been negative, and Washington has still asked Europe for a great deal. A patron can obtain the holder's surplus over exit, plus whatever it can threaten that does not depend on the holder's alignment. We state this for any rule that settles Europe's price at the end of a period, so that one result covers the bidding benchmark, bargaining under rupture risk, hold-up and price driving.

\subsection{Bite, budgets and assurance}\label{sec:slackprinciple}

Consider one period in which Europe is aligned and its premium is $Q$. Washington has finitely many demands. Demand $k$ costs Europe $t_k>0$ and is worth $w_k>0$ to the United States. A concession is \emph{reversible} if exit would undo it, so that it lowers Europe's aligned payoff and leaves its exit payoff unchanged. A concession is \emph{sunk} if it lowers both payoffs: payments, purchases, and transfers of assets or of control over them. \Cref{sec:linkage} shows that which concessions are reversible is itself a choice of Washington's.

\paragraph{Settlement.} At the end of the period Europe's price is settled at the premium then prevailing. Let $\mathcal P(q)$ be Europe's payoff from the settlement, relative to the aligned status quo, when its premium is $q$, and let $-\mathcal P(q)$ be Washington's. We require three properties: (S1) $\mathcal P(q)\ge q$, since Europe can always exit; (S2) $\mathcal P$ is nondecreasing; and (S3) $\mathcal P(q')-\mathcal P(q)\le q'-q$ for $q\le q'$, so that a higher premium never raises Europe's price by more than itself. Four settlements in this paper satisfy them while Washington retains Europe: the bidding payment $\mathcal P(q)=\pos{q}$ (\cref{lem:bidding}); bargaining under rupture risk, $\mathcal P(q)=\max\{q,\psi(q+\beta(V-q))\}$ (\cref{prop:master}), a maximum of affine functions with slopes $1$ and $\psi(1-\beta)$; hold-up with constant disagreement costs, which adds $(1-\psi)N$ to the second term (\cref{prop:pain}); and price driving, $\mathcal P(q)=\pos{q+\bar b_C-Q}$ (\cref{rem:spite}). Europe's \emph{surplus over exit} is $\Lambda_{\mathcal P}=\mathcal P(Q)-Q\ge 0$. In the bidding benchmark it is the slack $\Lambda=\pos{-Q}$.

\paragraph{Instruments.} Washington backs its demands with instruments (\cref{sec:security}). Withdrawing a set $J$ of instruments from an aligned Europe for the period lowers Europe's aligned payoff by $\Delta^A(J)$ and its exit payoff by $\Delta^X(J)$, with $0\le\Delta^X(J)\le\Delta^A(J)$ and both nondecreasing in $J$. We call
\[
u(J)=\Delta^X(J)\qquad\text{and}\qquad d(J)=\Delta^A(J)-\Delta^X(J)
\]
the unconditional and conditional parts of $J$. The conditional part is what Europe would lose by exiting anyway (for a single instrument, the $d_j$ of \cref{eq:Q}). The unconditional part is what Washington would still provide to an exited Europe and can threaten to take away. We impose no additivity: the AI instrument and US force posture may be complements or substitutes in Europe's payoff.

\paragraph{Ultimatums.} Washington's credible set $J^c$ is fixed at the start of the period, and no instrument outside it can be withdrawn. Washington issues \emph{adaptive ultimatums}: any sequence of demands, each naming unconceded demands and a threat set of instruments not yet withdrawn, each chosen after seeing Europe's earlier responses. If Europe refuses, Washington executes the threat at direct cost $e(J)\ge 0$, and executed instruments stay withdrawn for the period, so a threat once carried out is \emph{spent}. A threat that deters is not spent and can back the next demand. If $J_w$ has been withdrawn and Europe has made reversible concessions $T_{\mathrm r}$ and sunk ones $T_{\mathrm s}$, the settlement is played at the premium $Q+d(J_w)+T_{\mathrm r}$, and Europe receives $-\Delta^A(J_w)-T_{\mathrm r}-T_{\mathrm s}+\mathcal P(Q+d(J_w)+T_{\mathrm r})$. For a set $S$ of concessions, $T_{\mathrm r}(S)$ and $T_{\mathrm s}(S)$ are its reversible and sunk totals and $W(S)=\sum_{k\in S}w_k$ its value to Washington. Executing $J$ costs Washington $e(J)$ plus the rise in the settlement, which by (S3) is at most $d(J)$. After any history, executing a further set raises the premium by at most $\Delta^A(J^c)$, since $0\le d\le\Delta^A$. So if $e$ is nondecreasing, every subset of $J^c$ stays credible after every history when $\bar m_U\ge e(J^c)+\Delta^A(J^c)$, a condition the attainment claims below use and the bounds do not.

\begin{proposition}[Slack principle]\label{prop:slack}
Let the settlement $\mathcal P$ satisfy (S1)--(S3), let $\Lambda_{\mathcal P}=\mathcal P(Q)-Q$, and let
\[
B(J)=\Delta^A(J)-\big[\mathcal P\big(Q+d(J)\big)-\mathcal P(Q)\big],\qquad \varrho(T)=T-\big[\mathcal P(Q+T)-\mathcal P(Q)\big].
\]
No instrument outside $J^c$ can be withdrawn during the period, executing a withdrawal of $J$ costs the United States $e(J)\ge 0$ directly, and the United States retains Europe after every history.
\begin{enumerate}[label=(\alph*)]
\item \emph{Bite.} Refusing an ultimatum backed by $J$, issued when nothing has been conceded or withdrawn, costs Europe $B(J)$, net of the change in its settlement. $B$ is nondecreasing in $J$, and
\[ B(J)\le u(J)+\min\{d(J),\Lambda_{\mathcal P}\}, \]
with equality if the withdrawal holds Europe to its outside option ($\mathcal P(Q+d(J))=Q+d(J)$) or leaves its settlement unchanged ($\mathcal P(Q+d(J))=\mathcal P(Q)$). In the bidding benchmark equality always holds, $B(J)=u(J)+\min\{d(J),\Lambda\}$, and if $Q\ge 0$ only unconditional parts bite.
\item \emph{Reversible concessions.} Europe's net loss from reversible concessions totalling $T$ is $\varrho(T)\le\min\{T,\Lambda_{\mathcal P}\}$, with equality in the bidding benchmark. Faced with a single ultimatum that asks for reversible concessions $T_{\mathrm r}$ and sunk concessions $T_{\mathrm s}$ and is backed by $J$, Europe accepts if and only if $\varrho(T_{\mathrm r})+T_{\mathrm s}\le B(J)$.
\item \emph{One budget.}
  \begin{enumerate}[label=(\roman*)]
  \item In every subgame-perfect equilibrium of every adaptive sequence of ultimatums, Europe's total net loss is at most $B(J^c)$. Conditional instruments share one budget: $B(J)-u(J)\le\Lambda_{\mathcal P}$ for every $J$, and reversible concessions cost Europe at most $\Lambda_{\mathcal P}$ in total, however they are backed.
  \item The United States' equilibrium payoff is at most
  $V_1=\max\{W(S)-\mathcal P(Q+T_{\mathrm r}(S)):\ \varrho(T_{\mathrm r}(S))+T_{\mathrm s}(S)\le B(J^c)\}$, the payoff of the best single ultimatum backed by $J^c$.
  \item If every subset of $J^c$ stays credible and the United States can commit to a single ultimatum, one bundled ultimatum backed by $J^c$ yields $V_1$. Europe's loss then equals $B(J^c)$ if the United States has a divisible sunk demand of positive value and of size at least $B(J^c)$.
  \end{enumerate}
\item \emph{Assurance.} Without that commitment the United States still obtains $V_1$ if all demands are reversible and $B(J^c)\ge\Lambda_{\mathcal P}$, or if some best single ultimatum contains all of its demands. (In every equilibrium if $e(J)>0$ for every nonempty $J$, as with additive execution costs that are positive for each instrument; in some equilibrium otherwise.) In general it can obtain strictly less. Threats that deter are not spent, so the United States reuses them after Europe concedes, and Europe anticipates this. With one instrument with $d=0$ and $n$ identical sunk demands of cost $\tau$, the United States obtains exactly $n \bmod (b+1)$ of them, where $b=\lfloor B(J^c)/\tau\rfloor$.
\end{enumerate}
\end{proposition}

\Cref{app:proofs} has the proof and \cref{app:numerics} the brute-force checks. The proof uses only (S1)--(S3). Withdrawing a conditional part from an aligned Europe takes away something Europe would lose by exiting anyway, so it raises Europe's premium by $d(J)$. In the bidding benchmark the withdrawal hurts Europe up to $\Lambda$. Beyond $\Lambda$ the premium turns positive and Washington must pay rent to keep Europe aligned, which hands the excess back. A reversible concession works the same way, which is part (b): it lowers only the aligned payoff, so past $\Lambda$ Europe would exit unless compensated. Sunk concessions and unconditional parts escape the cap because they hurt Europe whatever it does.

The same holds whatever settles the price. Under rupture risk, agenda power can lift Europe's price above $\pos{Q}$, and the budget becomes $\Lambda_{\mathcal P}=P^*(Q)-Q$, which grows with $V$ while agenda power binds. It can also be smaller than $\Lambda$: when rupture risk makes Europe concede in the bargaining ($P^*(Q)<0$), threats can extract only the rest of its surplus over exit. Take $Q=-1$, $V=4$, $\beta=0.5$, $\lambda=1$, $r=0.25$ and $\phi=0$. Then $\Lambda=1$ and $\Lambda_{\mathcal P}=2.2$, and withdrawing a conditional part of $3$ costs Europe $1.8$. Whatever agenda power, hold-up or a rival's price driving wins Europe above its outside option, credible conditional threats can claw back, down to the exit premium. The agenda channel of \cref{cor:master}(e) therefore survives only to the extent that Washington lacks credible conditional threats large enough to push Europe to its outside option. Where the second term of $P^*$ binds, the settlement also rises at rate $\psi(1-\beta)$ with the premium, so a reversible concession of $T$ inside the budget costs Europe only $(1-\psi(1-\beta))T$: rupture risk returns part of every reversible concession.

\begin{figure}[t]
\centering
\includegraphics[width=0.95\textwidth]{figures/fig_slack.pdf}
\caption{The slack budget in the bidding benchmark. Panel (a): the bite of a conditional instrument with $d=2$ rises one for one with Europe's slack $\Lambda$ until $\Lambda=2$; an unconditional instrument with $u=1$ bites the same at every $\Lambda$. Panel (b): two conditional instruments with $d=1.5$ each. The sum of their separate bites (dashed) overstates their joint bite $\min\{3,\Lambda\}$ at every $\Lambda<3$. Bites from \cref{prop:slack}(a).}
\label{fig:slack}
\end{figure}

In panel (a) of \cref{fig:slack} the unconditional threat is the smaller one, yet it is the only one that still works as $\Lambda$ falls toward zero. Panel (b) illustrates part (c)(i). A coercer that prices each conditional threat on its own counts the same slack twice: at $\Lambda=1$ the separate bites sum to $2$ and the joint bite is $1$. The bound in (c)(i) needs only a fixed credible set (\cref{app:numerics}). Attaining it needs that set to stay credible after every history, which the cap condition stated before the proposition ensures.

Part (d) is the assurance problem of a coercer that cannot commit to stop demanding. Europe concedes only if everything the unspent threat will extract later stays within the bite. With $b+1$ identical sunk demands, a single ultimatum obtains $b$ of them. The adaptive sequence obtains none. A threat meant to induce compliance works only if the target believes compliance will end it \citep{schelling1966arms}, and \citet{pauly2024damned} shows that the steps that make a threat credible weaken that belief. A staked promise to make no further demand in the period buys the commitment. Stakes cost nothing. Part (d) therefore describes a patron whose cap $\bar m_U$ is small relative to the value of its later demands. Such a patron can gain by issuing an ultimatum it expects Europe to refuse, which burns an instrument that would otherwise back later demands (\cref{app:numerics}).

\begin{corollary}[Rationing]\label{cor:rationing}
In the bidding benchmark, suppose all demands are reversible and either the United States can commit to a single ultimatum or $B(J^c)\ge\Lambda$. Then its optimal demand set solves
\[ \max_{y_k\in\{0,1\}}\ \sum_k w_k y_k-\pos{\textstyle\sum_k t_k y_k-\Lambda}
   \quad\text{s.t.}\quad \min\Big\{\sum_k t_k y_k,\Lambda\Big\}\le B(J^c), \]
with one $y_k$ for each demand $k$. This is the knapsack $\max\sum_k w_k y_k$ s.t.\ $\sum_k t_k y_k\le\min\{\Lambda,B(J^c)\}$ if $B(J^c)<\Lambda$, or if demands are divisible ($y_k\in[0,1]$) and $w_k\le t_k$ for every $k$. In general the United States may exceed $\Lambda$ and pay the excess back as rent. The free capacity $\min\{\Lambda,B(J^c)\}$ falls as $Q$ rises toward $0$. Demands from every issue area draw on it.
\end{corollary}

With a general settlement the objective is $\sum_k w_ky_k-[\mathcal P(Q+\sum_kt_ky_k)-\mathcal P(Q)]$, and capacity is free only where $\mathcal P$ is flat.

\subsection{Linkage}\label{sec:linkage}

Whether a concession is reversible depends on what enforces it. In \cref{prop:slack} the only thing that holds Europe to a reversible concession is the threat of losing the conditional parts of Washington's instruments, so a Europe that exited would undo every such concession at once. That is why they share one budget. Washington can also enforce a concession with a penalty of its own, attached to the issue and not to Europe's alignment: a tariff threat behind a trade concession, or a sanction behind a purchase commitment.

Let demand $k$ carry an \emph{issue-specific penalty} $p_k\in[0,\infty]$: a staked threat, credible whatever Europe's alignment, to impose a loss $p_k$ on Europe if it reverses concession $k$. A payment that cannot be recovered has $p_k=\infty$. Washington \emph{links} a set $L$ of demands: it treats the reversal of a linked concession as exit, so reversing one while staying aligned gives Europe no more than exiting would. At the end of the period Europe chooses its alignment and which concessions to reverse.

\begin{proposition}[Linkage]\label{prop:linkage}
Maintain the assumptions of \cref{prop:slack}, with the penalties $p_k$ and the linked set $L$ as above.
\begin{enumerate}[label=(\alph*)]
\item \emph{Enforced and linked parts.} An aligned Europe delivers concession $k$ if and only if $k$ is linked or $t_k\le p_k$. Its \emph{enforced part} $\min\{t_k,p_k\}$ enters Europe's payoffs as a sunk concession, and its \emph{linked part} $\pos{t_k-p_k}$, when $k\in L$, as a reversible one. \Cref{prop:slack} holds with $T_{\mathrm s}(S)=\sum_{k\in S}\min\{t_k,p_k\}$ and $T_{\mathrm r}(S)=\sum_{k\in S\cap L}\pos{t_k-p_k}$.
\item \emph{Where demands crowd each other out.} Demands draw on the shared budget $\Lambda_{\mathcal P}$ only through their linked parts. If $t_k\le p_k$ for every demand, no demand needs linkage, accepted demands leave the settlement unchanged, and demands interact only through the bite constraint $\sum_{k\in S}t_k\le B(J^c)$. Linking a demand with $t_k\le p_k$ changes nothing.
\item \emph{Who gains from separate enforcement.} For a given set of concessions, raising $p_k$ weakly raises Europe's loss and weakly lowers the settlement Washington pays. If $B(J^c)\ge\sum_k t_k$, the payoff $V_1$ of Washington's best single ultimatum is nondecreasing in each $p_k$.
\end{enumerate}
\end{proposition}

The one budget of \cref{prop:slack} is thus the case in which Washington's only enforcement of reversible demands is the alliance itself ($p_k=0$). The more Washington can enforce a concession with a penalty that does not run through Europe's alignment, the less that concession draws on Europe's slack, and the more it costs Europe, because exit no longer undoes it. Europe's protection against an accumulation of demands is that they can be undone together. Crowding out through the slack should therefore appear between demands enforced only by linkage, and not between demands that Washington enforces separately, which interact only through the bite of its threats. That makes the cross-issue budget a claim about how Washington enforces its demands, which can be observed.

Export-control alignment and limits on servicing in China are enforced by linkage, beyond what flow jurisdiction blocks directly (\cref{sec:background}). The trade terms of August 2025, under which the EU said it would eliminate tariffs on US industrial goods and align its technology-security requirements with US ones \citep{whitehouse2025jointstatement}, are enforced at least partly by the threat of higher US tariffs, so only the part of them that a tariff threat cannot hold draws on Europe's slack. The NATO spending target for 2035 \citep{nato2025hague} and the EU's stated intention to buy more US military equipment \citep{whitehouse2025jointstatement} are closer to sunk, since they cost Europe whatever it later does with ASML. They draw on the bite of Washington's instruments and not on the slack. At The Hague in June 2025 Trump tied Spain's trade terms to its defence spending: ``We're negotiating with Spain on a trade deal. We're going to make them pay twice as much'' \citep{rte2025nato}. In our terms that is linkage, which Washington needs where an issue-specific penalty cannot hold a concession on its own.

\subsection{Counter-threats and hedges}\label{sec:counter}

As in \cref{sec:security}, Chinese coercion enters only as counter-threats attached to particular concessions, announced before Washington makes its demands. We also consider a hedge, legacy-tool sales and servicing worth $y_H$ to Europe and $w_H$ to Washington if stopped: Washington can demand that Europe stop it, backed by $J^c$ (linkage), or threaten to block the hedged flows with its own jurisdiction, at net ex post cost $e_F$.

\begin{proposition}[Counter-threats and hedges]\label{prop:counter}
Maintain the assumptions of \cref{prop:slack} in the bidding benchmark.
\begin{enumerate}[label=(\alph*)]
\item \emph{Counter-threats.} Before the United States moves, China chooses credible counter-threats $r_k\in[0,\bar r_k]$. If concession $k$ is made and $r_k>0$, China executes at cost $e^C_k>0$ and Europe bears the sunk loss $r_k$. China loses $z_k\ge 0$ when concession $k$ is made. The United States makes a single ultimatum backed by $J^c$, and its optimal bundle is unique at the profiles compared below.
  \begin{enumerate}[label=(\roman*)]
  \item A counter-threat adds $r_k$ to the sunk cost of $k$: Europe accepts $S$ if and only if
  \[ \min\{T_{\mathrm r}(S),\Lambda\}+T_{\mathrm s}(S)+\textstyle\sum_{k\in S}r_k\le B(J^c). \]
  \item China does better than with no threats only if the United States' demand set changes.
  \item If a counter-threat on $k$ is executed in China's optimal profile, then dropping it changes the United States' demand set to one that contains a demand outside the equilibrium set. China executes on path only to crowd out other demands, and this occurs for an open set of parameters.
  \item The United States' payoff is weakly lower under any profile of counter-threats.
  \item Europe's net loss is at most $B(J^c)$ with or without counter-threats. It equals $B(J^c)$ in both cases if the United States has a divisible sunk demand of positive value that China cannot threaten. The counter-threat then taxes the patron, not the holder. Otherwise Europe's loss can rise or fall.
  \end{enumerate}
\item \emph{Policing a hedge.} Stopping a hedge worth $y_H$ to Europe is a reversible concession. If it is the only demand, linkage stops the hedge if and only if $\min\{y_H,\Lambda\}\le B(J^c)$ and $w_H\ge\pos{y_H-\Lambda}$, where $w_H$ is its value to the United States. The United States then pays rent $\pos{y_H-\Lambda}$. Linkage is free only if $y_H\le\min\{\Lambda,B(J^c)\}$, and its free part $\min\{y_H,\Lambda\}$ vanishes as $Q\to 0$. A staked threat to block the hedge by flow jurisdiction, with net ex post cost $e_F$, is credible if and only if $e_F\le\bar m_U$. It leaves $Q$ unchanged and is never executed. Executing a conditional withdrawal $J$ raises Europe's premium to $Q+d(J)$.
\end{enumerate}
\end{proposition}

Part (a) changes who pays for Chinese coercion. By (a)(i) a counter-threat makes a concession more expensive for Europe exactly as a sunk cost would. If Washington has other demands that China cannot reach, it fills the budget with them, and Europe's loss stays at $B(J^c)$. Washington drops the demand China has made expensive for others it values less, so the counter-threat taxes the patron. China gains only by changing what Washington asks for (part (a)(ii)), and it carries out a threat only to crowd out another demand (part (a)(iii); \cref{app:proofs} gives an example on the equilibrium path). \Cref{sec:record} reads the Nexperia episode against that requirement.

Part (b) separates the slack logic from a coercer that prices each demand at a fixed cost of the cheapest credible instrument (\cref{sec:predictions}). The free part of linkage, $\min\{y_H,\Lambda\}$, shrinks as the premium rises toward zero, and beyond it Washington must pay rent. Flow jurisdiction does not depend on slack: it is credible whenever $e_F\le\bar m_U$ and leaves $Q$ unchanged. So linkage against a hedge worth less to Washington than to Europe ($w_H<y_H$) should appear when the ally has slack and fade as its premium approaches zero, while the use of flow jurisdiction should not track the premium at all. A cost minimizer that counted the rent as a cost of linkage would predict the same switch.

Given our reading that Europe's premium has been negative since 2019 (\cref{sec:evidence}), the record is consistent with the first half of this prediction. The 2025 National Security Strategy and the August 2025 joint statement tie favourable treatment and US AI exports to export-control alignment (\cref{sec:background}), a reversible concession. The second half needs variation in premia that the Dutch record alone does not supply (\cref{sec:tests}).

\subsection{Refusals}\label{sec:refusals}

Under complete information no demand is refused and no threat is carried out: Washington never makes a demand that Europe would resist. Yet Europe has resisted demands, such as limits on servicing in China (\cref{sec:intro}). A refusal needs Washington to be unsure what a demand would cost Europe. Let the net cost to Europe of one demand, $\kappa$, be Europe's private information, drawn from a continuous distribution $F$; its sign is unrestricted, since Europe may want some steps for its own reasons. To keep the settlement fixed, we let the demand be enforced ($t\le p$ in \cref{prop:linkage}), so that accepting it lowers Europe's aligned and exit payoffs equally. Washington values the concession at $w>0$, chooses a threat set $J\subseteq J^c$ before Europe responds, and carries out a refused threat. $J=\emptyset$ is an unbacked request, with $B(\emptyset)=0$. Executing $J$ costs Washington $E(J)=e(J)+\mathcal P(Q+d(J))-\mathcal P(Q)\ge 0$.

\begin{proposition}[Refusals and executed threats]\label{prop:refusal}
Under these assumptions Europe accepts if and only if $\kappa\le B(J)$, and the United States chooses $J$ to maximize $F(B(J))\,w-\big(1-F(B(J))\big)E(J)$.
\begin{enumerate}[label=(\alph*)]
\item In equilibrium the demand is refused with probability $1-F(B(J^*))$, and a refused demand that was backed ($J^*\neq\emptyset$) is followed by the execution of its threat. No threat set holds refusals below $1-F(B(J^c))\ge 1-F\big(u(J^c)+\min\{d(J^c),\Lambda_{\mathcal P}\}\big)$. With purely conditional instruments in the bidding benchmark this floor is at least $1-F(\Lambda)$, which rises as $Q$ rises toward $0$.
\item An unbacked request is refused with probability $1-F(0)$ and costs nothing when refused. Whenever $F(0)>0$, the United States prefers an unbacked request to no demand, so it always makes some demand, and if the demand it makes is unbacked, a refusal has no consequence.
\item Selecting the optimal threat set with the largest bite, $B(J^*)$ is nondecreasing in $w$: more valuable demands are backed by larger threats and are refused less often.
\item A refused and executed conditional threat raises Europe's premium by $d(J^*)$ and is spent for the rest of the period.
\end{enumerate}
\end{proposition}

The proposition says what the record should show. Near a zero premium, conditional threats bite little, so the refusals that Washington can prevent at any price become fewer (part (a)); the equilibrium refusal rate need not rise, since Washington may then switch to unconditional instruments or stop backing its demands. Requests that Washington does not back can be refused at no cost (part (b)). And the most valuable demands should come with the largest threats (part (c)). Dutch and Japanese resistance to limits on servicing in China fits an unbacked request better than a refused threat if, as the sources we use suggest, no US withdrawal followed it. The MATCH Act would back the request with a statutory deadline, but it had not been enacted by late October 2026 (\cref{sec:background}).

\subsection{Conditionality as a choice}\label{sec:conditionality}

Washington chooses its conditionality. With $\kappa=G(s_A)-G(s_X-m_S)$ from \cref{eq:kappa}, choosing the stake $m_S\le\bar m_U$ is choosing $\kappa\in[\kappa_0,\bar\kappa]$, where $\kappa_0=G(s_A)-G(s_X)$ is the conditionality with no stake and $\bar\kappa$ is the largest Washington can create. Raising $\kappa$ lowers Europe's premium by $S_k$ per unit. If Washington converts slack into concessions worth $\nu\in[0,1]$ per unit of European cost, and conditionality has a convex cost $H(\kappa)$ in allied assurance, \cref{prop:conditionality} in \cref{app:conditionality} characterizes its choice. A patron with no further demands ($\nu=0$) makes its guarantee just conditional enough to hold the premium at zero, or less if assurance becomes too costly first. Conditionality beyond that point is worth buying only because slack can be spent, so only a patron that plans to make demands makes its guarantee more conditional than the rivalry with China requires. When assurance is cheap at the margin, Washington answers a stronger Chinese offer by making its guarantee more conditional, by exactly enough to keep the premium at zero, and pays nothing. We read the favourable treatment that the National Security Strategy offers aligning countries this way (\cref{sec:counter,sec:coding}).

The assurance cost need not be assumed. Allies that doubt the patron build capabilities of their own \citep{lanoszka2018atomic}, and a guarantee tied to Europe's conduct can be triggered by mistake, so a second ally values its own alignment less when it sees Washington's guarantee to Europe made conditional.

\begin{remark}[A source of assurance cost]\label{rem:assurance}
Let a second ally's premium rise with Europe's conditionality, $Q_2(\kappa)=Q_2^0+a(\kappa-\kappa_0)$ with $a\ge 0$ and $Q_2^0\le 0$, and let Washington convert that ally's slack at the same rate $\nu$. Then $H(\kappa)=\Pi_2(\kappa_0)-\Pi_2(\kappa)$, with $\Pi_2=-\pos{Q_2}+\nu\pos{-Q_2}$, is nondecreasing and convex with $H(\kappa_0)=0$, so \cref{prop:conditionality} applies. Let $\tilde\kappa=\kappa_0-Q_2^0/a$ be the conditionality at which the second ally's premium reaches zero. If Europe's premium is negative at $\kappa_0$ and $\nu>0$, an optimal conditionality is $\kappa^*=\kappa_0$ if $a\ge S_k$, $\kappa^*=\min\{\tilde\kappa,\bar\kappa\}$ if $\nu S_k\le a<S_k$, and $\kappa^*=\bar\kappa$ if $a<\nu S_k$ (with $\tilde\kappa=\infty$ if $a=0$).
\end{remark}

Conditionality toward one ally spends the slack of the others. Washington makes Europe's guarantee more conditional only if Europe's stake in it, $S_k$, exceeds the spillover $a$ onto the second ally, and it stops where that ally's premium reaches zero unless the spillover is small even relative to the value of Europe's concessions. A patron with many allies watching makes its guarantees less conditional than one bargaining with a single holder.

An untargeted drawdown, such as the cut of 25,000 or more troops from Europe that NBC News reported in September 2026 the Pentagon was weighing \citep{nbc2026pentagon}, lowers what Washington provides to a loyal ally, $x_p(A)$, and with it $d_p$ and Europe's slack. The May 2026 order moving troops out of Germany (\cref{sec:background}) punished an act with no prior demand, which lies outside our demand-and-threat structure. Read as the execution of an implicit threat, it spends the instrument (\cref{sec:slackprinciple}) and raises the target's premium by the instrument's conditional part (\cref{prop:counter}(b)). Carrying it out at all requires the burning logic described after \cref{prop:slack}, or incomplete information as in \cref{prop:refusal}, and its effect on other allies' beliefs about $\kappa$ is what \cref{rem:assurance} prices.

The silicon shield, the fourth channel of \cref{cor:master}(e), has no fixed sign. If holding the chokepoint raises Washington's stake in Europe's security whatever Europe's alignment, $s_X$ rises. Where $G$ is unimodal with its mode at or below $s_X-\bar m_U$, so that the most likely crisis costs lie below Washington's net stake in a Europe that has exited, a higher $s_X$ lowers the most conditionality Washington can create (\cref{prop:conditionality}(d)). Where a staked abandonment makes defence of an exiter unlikely, $s_X-\bar m_U$ lies in the left tail of $G$, where the density rises, and holding the chokepoint can make the guarantee more conditional. \citet{blankenship2021price} finds that allies in locations the United States values contribute less to the common burden, which is what the first case implies.
